\documentclass[12pt, spanish]{article}
\usepackage[utf8]{inputenc}
\usepackage[T1]{fontenc}
\usepackage[spanish]{babel}
\usepackage{amsmath, amssymb}
\usepackage{geometry}
\usepackage{enumitem}
\usepackage{xcolor}
\usepackage{mdframed}

\geometry{margin=2.5cm}

\definecolor{correctcolor}{RGB}{0, 120, 0}
\definecolor{headercolor}{RGB}{30, 60, 120}

\newcounter{pregunta}
\newenvironment{preguntabox}[1]{%
  \stepcounter{pregunta}
  \begin{mdframed}[linecolor=headercolor, linewidth=1.2pt, topline=true, bottomline=true]
  \noindent{\color{headercolor}\textbf{Pregunta \thepregunta\ — #1}}\\[6pt]
}{%
  \end{mdframed}\vspace{8pt}
}

\newcommand{\correcta}[1]{\item[\textcolor{correctcolor}{\textbf{(#1)}}] \textcolor{correctcolor}}
\newcommand{\opcion}[2]{\item[(#1)] #2}

\title{\textbf{Preguntas de Mayor Dificultad}\\[4pt]
       \large Control de Procesos — Análisis de Sistemas}
\author{Universidad ECCI}
\date{}

\begin{document}
\maketitle
\thispagestyle{empty}

\begin{center}
\textit{Las 5 preguntas más desafiantes del banco, una por tema.\\
La respuesta correcta aparece resaltada en verde al final de cada pregunta.}
\end{center}

\vspace{10pt}
\hrule
\vspace{12pt}

% ─── PREGUNTA 1 ───────────────────────────────────────────────────────────────
\begin{preguntabox}{Conceptos de EDOs}
La solución completa de un sistema ante una entrada escalón es:
\[
  y(t) = 5 + e^{-t}\bigl[A\cos(3t) + B\sin(3t)\bigr]
\]
¿Qué revela esta expresión sobre el sistema?

\begin{enumerate}[label=(\Alph*), leftmargin=2em, itemsep=4pt]
  \correcta{A}{Es de segundo orden oscilatorio estable: la parte imaginaria $\pm 3j$ indica oscilación, la parte real $-1$ indica decaimiento, y la constante $5$ es el estado estacionario.}
  \opcion{B}{Es inestable porque contiene funciones trigonométricas que no convergen a un valor fijo con el tiempo.}
  \opcion{C}{El valor final es cero porque las exponenciales siempre dominan sobre las constantes cuando $t \to \infty$.}
  \opcion{D}{Es de primer orden porque tiene un solo coeficiente de amortiguamiento $(-1)$ en la exponencial.}
\end{enumerate}

\vspace{6pt}
\noindent\textbf{Explicación:} La estructura $e^{-t}[A\cos(3t)+B\sin(3t)]$ corresponde a raíces
$s = -1 \pm 3j$. La parte real negativa garantiza estabilidad; la parte imaginaria produce
oscilación con frecuencia $3\,\text{rad/s}$. Cuando $t\to\infty$ la solución homogénea decae
y solo permanece $y_p = 5$ (estado estacionario). El sistema es de segundo orden oscilatorio estable.
\end{preguntabox}

% ─── PREGUNTA 2 ───────────────────────────────────────────────────────────────
\begin{preguntabox}{Sistemas de Primer Orden}
De un proceso se sabe que su respuesta al escalón llega al $50\%$ del valor final en $t = 2\,\text{s}$.
¿Cuál es $\tau$?

\begin{enumerate}[label=(\Alph*), leftmargin=2em, itemsep=4pt]
  \correcta{A}{$\tau \approx 2{,}89\,\text{s}$, porque el $50\%$ se alcanza en
    $t_{50} = \tau\ln 2 \approx 0{,}693\,\tau$, despejando $\tau = \dfrac{2}{0{,}693}$.}
  \opcion{B}{$\tau = 2\,\text{s}$, porque el tiempo al $50\%$ siempre es igual a la constante de tiempo del sistema.}
  \opcion{C}{$\tau = 4\,\text{s}$, porque el tiempo al $50\%$ es siempre la mitad de la constante de tiempo real.}
  \opcion{D}{$\tau$ no puede determinarse: el tiempo al $50\%$ no está relacionado con $\tau$ de forma directa.}
\end{enumerate}

\vspace{6pt}
\noindent\textbf{Explicación:} En $t_{50}$:
\[
  0{,}5 = 1 - e^{-t_{50}/\tau}
  \;\Rightarrow\;
  e^{-t_{50}/\tau} = 0{,}5
  \;\Rightarrow\;
  t_{50} = \tau\ln 2 \approx 0{,}693\,\tau
  \;\Rightarrow\;
  \tau = \frac{2}{0{,}693} \approx 2{,}89\,\text{s}.
\]
El $50\%$ no ocurre en $t = \tau$ sino en $t \approx 0{,}693\,\tau$. Cada porcentaje de la
respuesta corresponde a un múltiplo distinto de $\tau$.
\end{preguntabox}

% ─── PREGUNTA 3 ───────────────────────────────────────────────────────────────
\begin{preguntabox}{Respuesta de Primer Orden}
Respuesta al escalón de un proceso: la salida permanece en cero durante $1\,\text{s}$, luego
sube como exponencial hasta el valor final sin sobreimpulso. El ajuste desde $t = 1\,\text{s}$
entrega $\tau = 4\,\text{s}$ y $K = 2$. ¿Cuál es el modelo más apropiado?

\begin{enumerate}[label=(\Alph*), leftmargin=2em, itemsep=4pt]
  \correcta{A}{Primer orden con tiempo muerto:
    $y(t) = K\,\Delta u\!\left(1 - e^{-(t-1)/4}\right)$ para $t > 1\,\text{s}$,
    con $\theta = 1\,\text{s}$, $\tau = 4\,\text{s}$, $K = 2$.}
  \opcion{B}{Primer orden puro con $\tau_{\text{efectivo}} = 5\,\text{s}$, porque el tiempo muerto
    se debe incluir en la constante de tiempo.}
  \opcion{C}{Segundo orden no oscilatorio, porque el retardo inicial indica la presencia de dos modos
    dinámicos distintos.}
  \opcion{D}{Primer orden con $\tau = 4\,\text{s}$ y $K = 2$, ignorando el tiempo plano porque es
    probablemente ruido de medición.}
\end{enumerate}

\vspace{6pt}
\noindent\textbf{Explicación:} El tiempo plano de $1\,\text{s}$ es el tiempo muerto $\theta$: el
sistema no responde en absoluto durante ese intervalo. El modelo correcto es primer orden con tiempo
muerto ($\theta = 1\,\text{s}$, $\tau = 4\,\text{s}$, $K = 2$). Incluir $\theta$ en $\tau$ produce
un modelo inexacto que no refleja la física real del proceso.
\end{preguntabox}

% ─── PREGUNTA 4 ───────────────────────────────────────────────────────────────
\begin{preguntabox}{Sistemas de Segundo Orden}
La forma estándar del sistema de segundo orden es:
\[
  \frac{1}{\omega_n^2}\frac{d^2y}{dt^2} + \frac{2\zeta}{\omega_n}\frac{dy}{dt} + y = K\,u
\]
¿Qué valor de $\zeta$ marca el límite entre respuesta oscilatoria y no oscilatoria?

\begin{enumerate}[label=(\Alph*), leftmargin=2em, itemsep=4pt]
  \correcta{A}{$\zeta = 1$: para $\zeta < 1$ la respuesta es oscilatoria (subamortiguada); para
    $\zeta > 1$ es no oscilatoria (sobreamortiguada).}
  \opcion{B}{$\zeta = 0{,}5$: ese valor separa los dos tipos de respuesta de forma óptima.}
  \opcion{C}{$\zeta = 0$: sin amortiguamiento el sistema siempre oscila, pero no hay límite entre
    los dos tipos.}
  \opcion{D}{$\zeta = 2$: ese es el coeficiente de la primera derivada en la forma estándar.}
\end{enumerate}

\vspace{6pt}
\noindent\textbf{Explicación:} Las raíces de la ecuación característica son:
\[
  s = -\zeta\omega_n \pm \omega_n\sqrt{\zeta^2 - 1}.
\]
Para $\zeta < 1$: raíces complejas conjugadas $\Rightarrow$ subamortiguado (oscilatorio).
Para $\zeta = 1$: raíz doble real $\Rightarrow$ críticamente amortiguado.
Para $\zeta > 1$: dos raíces reales distintas $\Rightarrow$ sobreamortiguado.
El comportamiento oscilatorio observado en clase corresponde a $\zeta < 1$.
\end{preguntabox}

% ─── PREGUNTA 5 ───────────────────────────────────────────────────────────────
\begin{preguntabox}{Respuestas Dinámicas}
Un proceso completamente desconocido debe identificarse con un solo experimento.
¿Cuál es la estrategia más informativa y práctica?

\begin{enumerate}[label=(\Alph*), leftmargin=2em, itemsep=4pt]
  \correcta{A}{Respuesta al escalón: da $K$ del valor final, $\tau$ del $63\%$, y la forma de la
    curva revela orden superior, tiempo muerto u oscilaciones.}
  \opcion{B}{Respuesta al impulso: contiene toda la dinámica del sistema y es el experimento más
    breve posible.}
  \opcion{C}{Respuesta a la rampa: el error de seguimiento estacionario da simultáneamente $K$ y
    $\tau$ en un solo dato.}
  \opcion{D}{Respuesta senoidal a múltiples frecuencias: es la más completa y la más rápida de
    realizar en planta.}
\end{enumerate}

\vspace{6pt}
\noindent\textbf{Explicación:} El escalón es el estándar industrial: fácil de aplicar, entrega $K$
del valor final, $\tau$ del $63\%$, y la forma de la curva diagnostica el tipo de modelo
(curva en S $\to$ orden superior; sobreimpulso $\to$ oscilatorio; retardo $\to$ tiempo muerto).
El impulso es completo en teoría pero imposible de aplicar con precisión en planta.
\end{preguntabox}

\end{document}
